\section{B. The Fermat tower and its homology}

We use the arithmetic model actually written in (220):
\[
 X_n^+:\ x_0^n+x_1^n+x_\infty^n=0,\qquad
 D_n=\{x_0x_1x_\infty=0\},\qquad U_n^+=X_n^+\setminus D_n.
\]
The notation \(X_n:x^n+y^n=z^n\), used in our geometric chain
calculations, is related to this model explicitly in B.1. Throughout this
part, \(F=\langle x,y\rangle\), \(c=[x,y]=xyx^{-1}y^{-1}\), and
\(F''=[F',F']\). In a profinite group, a commutator subgroup means its
closed version. Put
\[
 A_n=(\mathbb Z/n)^2,\quad
 R_n=\mathbb Z[A_n]=
 \mathbb Z[X^{\pm1},Y^{\pm1}]/(X^n-1,Y^n-1),
 \quad K_n=\ker(F\longrightarrow A_n).
\]
Capital letters \(X,Y\) in a group ring denote the images of \(x,y\).
We write
\(M=\mathbb Z^3/\mathbb Z(1,1,1)\), with basis the classes \(e_0,e_1\)
and \(e_\infty=-e_0-e_1\).

\begin{statement}{B.1}{Proved}{Arithmetic models and the precise problem}
The morphism
\[
 q_n:U_n^+\longrightarrow U_1^+,\qquad
 [x_0:x_1:x_\infty]\longmapsto[x_0^n:x_1^n:x_\infty^n]
\]
is defined over \(\mathbb Q\), is equivariant for coordinate permutations,
and has deck group \(\mu_n^3/\mu_n\). The coordinate
\[
 z=-x_0/x_\infty,\qquad
 \psi(z)=[-z:z-1:1]
\]
identifies \(U_1^+\) with \(\mathbb P^1-\{0,1,\infty\}\).
For \(\alpha_n^n=-1\), the geometric comparison
\[
 \iota_n:X_n^+\longrightarrow X_n,\qquad
 [x_0:x_1:x_\infty]\longmapsto[x_0:x_1:\alpha_n x_\infty]
\]
is an isomorphism over \(\mathbb Q(\alpha_n)\).
At odd levels one can take \(\alpha_n=-1\) over \(\mathbb Q\);
at level four the displayed comparison uses \(\alpha_4=\zeta_8\).
In particular \(X_4^+(\mathbb R)=\varnothing\), whereas
\([0:1:1]\in X_4(\mathbb Q)\).

The \emph{position du probl\`eme} in \S51.11 is to identify the affine
representation
\[
 \widetilde{\mathcal M}\longrightarrow
 G_{\rho,\sigma}/\mu_G
 \longrightarrow \widehat{\mathbb Z}^{\,*}\times S_3
\]
with the action on a fibre of the universal abelian Fermat covering,
and separately to interpret its specialization at \(z=1/2\).
Its translation group is \(\widehat M(1)\).
\emph{Manuscript: \S51.11, folios 695--697, scans 95--99, (218)--(230);
\S51.14, folio 723, scan 151, (315)--(317).}
\end{statement}
\begin{proof}
The asserted equations follow on taking \(n\)-th powers. On the chart
\(x_\infty\ne0\), division by \(\alpha_n\) gives
\((x_0/\alpha_nx_\infty)^n+(x_1/\alpha_nx_\infty)^n=1\);
the induced coordinate on the base is \(z\). Compatible choices, for
example \(\alpha_n=\exp(\pi i/n)\) after fixing an embedding in
\(\mathbb C\), commute with the power maps. Thus all later geometric
comparisons have an actual map behind them. The even-level real assertion
follows because a sum of three real fourth powers vanishes only when
all three coordinates vanish. Equations (228)--(230) explicitly identify
the deck lattice and its cyclotomic action; (219) singles out the remaining
rank-one affine specialization. These are the data just stated.
\end{proof}

\begin{statement}{B.2}{Proved}{The universal abelian covering}
Over \(\mathbb C\), \(q_n\) corresponds to the subgroup \(K_n\) of
\(\pi_1(U_1)=F\). Its monodromy is \(x\mapsto(1,0)\),
\(y\mapsto(0,1)\). Consequently the pro-\'etale tower
\(\varprojlim_n U_n^+\) is the universal abelian geometric covering, with
deck module \(\widehat M(1)\). Its deck and coordinate-permutation
actions form \((\mu_n^3/\mu_n)\rtimes S_3\) at finite level.
\emph{Manuscript: \S51.11, folios 696--697, scans 97--99, (220)--(230);
\S51.13, folio 715, scan 135.}
\end{statement}
\begin{proof}
In the standard geometric chart, \(u^n=z\) and \(v^n=1-z\). Analytic
continuation around \(0\) multiplies \(u\) by \(\zeta_n\) and leaves \(v\)
unchanged; continuation around \(1\) does the reverse. These two elements
generate \(A_n\), so the covering is connected and its subgroup is exactly
\(K_n\). Every finite quotient of \(\widehat{\mathbb Z}^2\) factors through
some \(A_n\); this proves the universal assertion, rather than only the
existence of a collection of abelian coverings. A change of compatible
geometric roots identifies the deck group with \(\widehat{\mathbb Z}^2\)
by a change of basis; intrinsically it is
\((\widehat{\mathbb Z}^3/\widehat{\mathbb Z})\otimes
\widehat{\mathbb Z}(1)\). Permutations act on the three coordinates,
commute with power maps, and conjugate the diagonal multipliers by
permuting them. This proves the stated semidirect product and the
compatibility of the action with the inverse system.
\end{proof}

\begin{statement}{B.3}{Refuted}{The withdrawn freeness and its exact replacement}
For \(n>1\), the action of \(\mu_n^3\) on \(X_n^+\) is not free:
its diagonal \(\mu_n\) fixes the whole curve. The quotient group scheme
\(A_n^{\rm geom}=\mu_n^3/\mu_n\) acts freely on \(U_n^+\), and
\(U_n^+\to U_1^+\) is an \(A_n^{\rm geom}\)-torsor.
At a coordinate cusp its stabilizer in \(A_n^{\rm geom}\) is a copy of
\(\mu_n\). Thus the quotient action on the compact curve is not free.
These are assertions about actions on schemes; the corresponding
integral homology modules are computed in B.5--B.7.
\emph{Manuscript: \S51.11, folio \(696'\), scan 98, deletion before (224)
and the inserted \(U_n\) in the paragraph following (225).}
\end{statement}
\begin{proof}
A multiplier \((a_0,a_1,a_\infty)\) fixes a point with all coordinates
invertible precisely when the three multipliers agree projectively.
Its class in the quotient is then the identity. At a point
\([0:b:c]\), with \(bc\ne0\), the equality required is
\(a_1=a_\infty\); \(a_0\) is unrestricted. Modulo the diagonal this is
exactly \(\mu_n\). This calculation also describes the stabilizer
functors, and hence applies to group schemes, including their
nonreduced fibres.
On a torus chart the power map is given by adjoining independent
\(n\)-th roots of two units. Its coordinate algebra is locally free of
rank \(n^2\), and the standard root-of-unit torsor identity proves
faithful flatness and the torsor assertion. Over \(\mathbb Z[1/n]\)
it is finite \'etale. At residue characteristics dividing \(n\) it is
still a torsor for the fppf topology; this does not assert \'etaleness or
smoothness of those Fermat fibres.
The finite checks in \texttt{code/b\_freeness.py} enumerate the
full and quotient stabilizers over finite fields for \(1\le n\le8\).
They give full-group stabilizer order \(n\) on the open part,
quotient stabilizer order one there, and quotient stabilizer
order \(n\) at a cusp.
\end{proof}

\begin{statement}{B.4}{Proved}{The discrete Alexander module}
Conjugation induces a free rank-one
\(R=\mathbb Z[X^{\pm1},Y^{\pm1}]\)-module structure
\[
 F'/F''=R\,c.
\]
With the convention for \(c\) fixed above, the injective Fox-chain map is
\[
 R\,c\longrightarrow R^2,\qquad
 f c\longmapsto ((1-Y)f,(X-1)f),
\]
and its image is
\(\ker\bigl((X-1,\;Y-1):R^2\to R\bigr)\).
\end{statement}
\begin{proof}
The graph with one vertex and oriented edges \(x,y\) has fundamental
group \(F\). Its universal abelian covering has vertex module \(R\),
edge module \(R^2\), and the displayed boundary. Its first homology is
the abelianization of its fundamental group \(F'\). The chain of
\(xyx^{-1}y^{-1}\) is \((1-Y,X-1)\).
For a cycle \((a,b)\), reduction of
\((X-1)a+(Y-1)b=0\) modulo \(X-1\) gives
\((Y-1)\bar b=0\) in the domain \(\mathbb Z[Y^{\pm1}]\).
Thus \(b=(X-1)f\). Substitution and cancellation of the nonzero
divisor \(X-1\) give \(a=(1-Y)f\). The same cancellation proves
uniqueness. This supplies both the generator and the proof of freeness.
\end{proof}

\begin{statement}{B.5}{Proved}{All-cusp-deleted integral homology}
There are canonical \(R_n\)-linear descriptions
\[
 H_1(U_n,\mathbb Z)=
 \ker\left(R_n^2\xrightarrow{(X-1,\;Y-1)}R_n\right),             \tag{B.1}
\]
\[
 0\longrightarrow
 R_n/\mathbb Z\,\mathcal N_n
 \xrightarrow{\ f\mapsto f c\ }H_1(U_n,\mathbb Z)
 \xrightarrow{q_{n*}} n\mathbb Z^2\longrightarrow0,\qquad
 \mathcal N_n=\sum_{a,b=0}^{n-1}X^aY^b.                       \tag{B.2}
\]
The deck action on the right-hand term is trivial.
The group \(H_1(U_n,\mathbb Z)\) has rank \(n^2+1\), and over
\(\mathbb C\) the trivial deck character occurs twice and every other
character occurs once.
\end{statement}
\begin{proof}
The finite covering of the two-circle graph has \(n^2\) vertices and
\(2n^2\) edges, with boundary as in (B.1). A spanning tree proves that
its cycle group is free of rank \(2n^2-n^2+1\).
The inclusion \(K_n\subset F\) induces \(K_n^{\rm ab}\to F^{\rm ab}\).
Its image is \(n\mathbb Z^2\); its kernel is the image of \(F'\),
because \(F'\subset K_n\) and \(K_n/F'=n\mathbb Z^2\).
B.4 identifies that image with the orbit module generated by \(c\).
The kernel of \(R_n\to R_n^2,\ f\mapsto ((1-Y)f,(X-1)f)\)
consists of elements fixed by both translations. Equality of all
coefficients says precisely that it is \(\mathbb Z\mathcal N_n\).
This proves (B.2), including its exactness over \(\mathbb Z\).
On a deck character \((\zeta_n^a,\zeta_n^b)\), the boundary in
(B.1) has rank one unless \(a=b=0\); in that case it has rank zero.
This proves the character formula directly.
\end{proof}

\begin{statement}{B.6}{Proved}{The compact integral homology module}
Define
\[
 S_X=\sum_{i=0}^{n-1}X^i,\quad S_Y=\sum_{j=0}^{n-1}Y^j,\quad
 P_n=\sum_{0\le j\le i<n}X^iY^j,\quad
 I_n=(S_X,S_Y,P_n)\subset R_n.
\]
The compact homology has the exact integral presentation
\[
 H_1(X_n,\mathbb Z)\simeq R_n/I_n,\qquad 1\longmapsto c.        \tag{B.3}
\]
Equivalently, (B.1) is divided by the \(3n\) cusp cycles
\[
 Y^j(S_X,0),\quad X^i(0,S_Y),\quad
 X^i(S_{XY},X S_{XY})\quad(0\le i,j<n),
 \qquad S_{XY}=\sum_{i=0}^{n-1}(XY)^i.                        \tag{B.4}
\]
The cusp subgroup is primitive of rank \(3n-1\).
The rank of (B.3) is \((n-1)(n-2)\), and its complex characters
are exactly those \((a,b)\) with \(a,b,a+b\ne0\pmod n\), each once.
\end{statement}
\begin{proof}
Filling a puncture kills its oriented peripheral loop in first homology.
Lifting the peripheral words \(x^n,y^n,(xy)^n\) gives precisely (B.4).
The homology sequence of the pair consisting of the compact surface and
its punctured surface identifies the subgroup of peripheral classes
with \(\mathbb Z[D_n]/\mathbb Z\sum_{P\in D_n}P\).
The quotient is the integral homology of an oriented compact surface,
so is torsion-free. This proves primitivity and rank \(3n-1\).

We give an integral derivation of the ring presentation. First fill
the \(2n\) cusps over \(0,1\). The resulting affine curve
\(u^n+v^n=1\) retracts onto the inverse image of the segment \([0,1]\).
To see this, take a closed regular neighbourhood of that segment
in the \(z\)-plane. Outside it the finite map is unramified and its
ends are collars. Inside it the local charts at the only branch
points are \(w\mapsto w^n\); their lifted discs and the intervening
strips are a regular neighbourhood of the lifted segment.
Shrinking these collars, discs and strips gives the stated retraction.
That graph has vertices \(A_n/\langle X\rangle\) and
\(A_n/\langle Y\rangle\), and one edge for each element of \(A_n\);
it is the complete bipartite graph on two sets of \(n\) vertices.
Its cycles are the integral matrices whose row and column sums vanish.
The map
\[
 R_n\longrightarrow\{\text{such matrices}\},\qquad
 f\longmapsto(1-X)(1-Y)f
\]
is surjective: subtracting the last row and column expresses every such
matrix as an integral sum of elementary rectangles. Its kernel is
\((S_X,S_Y)\), as one sees by successively taking differences of rows
and columns. The elementary rectangle is the commutator cycle, up to
the fixed choice of its orientation.

Now fill the \(n\) remaining cusps. In the original two-edge chain
module the identities
\[
 (1-Y)P_n=S_X-Y S_{XY},\qquad
 (X-1)P_n=S_Y-S_{XY}
\]
give
\[
 (S_{XY},X S_{XY})
 =-Y^{-1}P_n\,c+Y^{-1}(S_X,0)+Y^{-1}(0,S_Y).
\]
All infinity cusp cycles are translates of this one. Thus their
classes in the bipartite graph kill exactly the ideal generated by
\(P_n\), proving (B.3).
For a character with both \(a,b\ne0\), the last two displayed identities
make \(P_n\) zero precisely when \(a+b\ne0\); when \(a+b=0\),
\((X-1)P_n=-n\) is nonzero. The character statement follows. Its
cardinality is \((n-1)(n-2)\), agreeing with the topological rank.
\end{proof}

\begin{statement}{B.7}{Proved}{Finite integral computations and module freeness}
The exact calculations for \(1\le n\le8\) give
\[
\begin{array}{c|rrrrrrrr}
n&1&2&3&4&5&6&7&8\\ \hline
g(X_n)&0&0&1&3&6&10&15&21\\
\operatorname{rk}H_1(U_n)&2&5&10&17&26&37&50&65\\
\operatorname{rk}H_1(X_n)&0&0&2&6&12&20&30&42\\
\operatorname{rk}\langle\text{cusps}\rangle&2&5&8&11&14&17&20&23
\end{array}
\]
For every \(n>1\), \(H_1(U_n,\mathbb Z)\) is not a free \(R_n\)-module.
For every \(n\ge3\), \(H_1(X_n,\mathbb Z)\) is cyclic but not free
as an \(R_n\)-module. At \(n=1,2\) the compact module is zero.
The exact comparison maps, deck actions and \(S_3\)-actions are supplied
as integer matrices in \texttt{code/b\_fermat\_modules.json}.
\end{statement}
\begin{proof}
The program builds the lifted graph, chooses the tree obtained from the
paths \(x^i y^j\), and uses each remaining edge as a fundamental cycle.
It checks the Fox boundary, computes the deck transformations and the
two based substitutions in B.9, and expresses every cusp cycle in
this integral basis. It checks that the nonzero Smith entries of the
cusp matrix and of the compact annihilator presentation are all one.
It also checks that the map to \(F^{\rm ab}\) has Smith entries \(n,n\).
These are full integer presentations, rather than character counts.
All output rows displayed here were generated and checked by that program.

A free \(R_n\)-module has integral rank divisible by \(n^2\).
The affine rank \(n^2+1\) fails this test for \(n>1\).
The positive compact rank \((n-1)(n-2)\) is strictly below \(n^2\).
These rank conclusions apply through the explicit presentations
(B.1) and (B.3); they do not replace the scheme-action assertion of B.3.
\end{proof}

\begin{statement}{B.8}{Proved}{The full profinite rank-one theorem}
Let \(\widehat F\) be the free profinite group on \(x,y\), and put
\[
 \Lambda=\widehat{\mathbb Z}[[\widehat{\mathbb Z}^{\,2}]]
 =\varprojlim_{m,n}(\mathbb Z/m)[A_n].
\]
There are continuous \(\Lambda\)-linear isomorphisms
\[
 \widehat F'/\widehat F''
 \xrightarrow{\ \sim\ }\varprojlim_n H_1(U_n,\widehat{\mathbb Z})
 \xrightarrow{\ \sim\ }\varprojlim_n H_1(X_n,\widehat{\mathbb Z})
 \simeq\Lambda\,c.                                          \tag{B.5}
\]
The transition maps on homology are pushforwards of the covering maps.
The same limits can be formed diagonally with coefficients
\(\mathbb Z/n\). This is the full profinite assertion, including
the prime-to-\(\ell\) factors of the completed group ring.
\end{statement}
\begin{proof}
We first prove the group-to-tower comparison. Let
\(\widehat K_n=\ker(\widehat F\to A_n)\) and
\(H=\bigcap_n\widehat K_n=\widehat F'\).
The map \(H\to\varprojlim_n\widehat K_n^{\rm ab}\) is surjective:
compatible abelianized cosets are nonempty compact subsets of
\(\widehat K_n\), nested along a cofinal divisibility chain; their
intersection supplies a lift in \(H\).
Its kernel contains \(H'\). To prove the reverse inclusion, take a
continuous homomorphism \(f:H\to C\) to a finite abelian group. Choose
an open normal subgroup \(V\) of \(\widehat F\) with
\(V\cap H\subset\ker f\). Since \(H\) is normal,
\([V,H]\subset V\cap H\), and \(f\) extends to
\(HV\) by \(f(hv)=f(h)\). Compactness gives
\(\widehat K_n\subset HV\) for a sufficiently divisible \(n\).
It follows that every element in
\(\bigcap_n\widehat K_n'\) is annihilated by \(f\). Finite abelian
quotients separate \(H/H'\), proving the first isomorphism.
The covering graph identifies
\(\widehat K_n^{\rm ab}\) with \(H_1(U_n,\widehat{\mathbb Z})\):
the finite-index topology on \(K_n\subset F\) is its full profinite
topology, because the intersection of the finitely many \(F\)-conjugates
of a finite-index subgroup of \(K_n\) is finite-index normal in \(F\).
The extra quotient \(n\widehat{\mathbb Z}^2\) in (B.2) has transition
maps given by inclusion and has inverse limit
\(\bigcap_n n\widehat{\mathbb Z}^2=0\). Thus its rank-two contribution
at every finite level does not survive in (B.5).

For the second isomorphism, work first modulo \(m\).
The kernel of \(H_1(U_n,\mathbb Z/m)\to H_1(X_n,\mathbb Z/m)\)
is the peripheral module of B.6. Under \(U_{dn}\to U_n\), a small
loop around a cusp maps to \(d\) times a small loop. Thus the image of
that kernel is zero when \(m\mid d\). The inverse system is pro-zero.
All groups are finite, so taking the inverse limit preserves the short
exact sequences. This gives the second isomorphism modulo \(m\),
and passage to the inverse limit over \(m\) gives the assertion.

Finally use (B.3) modulo \(m\). It stays exact because its integral
quotient is free. Under \(R_{dn}\to R_n\),
\[
 S_{X,dn}\mapsto dS_{X,n},\quad
 S_{Y,dn}\mapsto dS_{Y,n},\quad
 P_{dn}\mapsto dP_n+\binom d2\mathcal N_n.
\]
The last formula counts pairs \(0\le j\le i<dn\) with each prescribed
pair of residues modulo \(n\). A value of \(d\) divisible by \(2m\)
kills every image generator of \(I_{dn}\) modulo \(m\).
Hence \(\varprojlim_n I_n(\mathbb Z/m)=0\), and the finite inverse-limit
exact sequence identifies \(\varprojlim_n H_1(X_n,\mathbb Z/m)\)
with \(\varprojlim_n(\mathbb Z/m)[A_n]\), by the compatible generator
\(c\). Passing to \(m\) proves freeness. This argument agrees with
\cite[Theorem 6.1]{Otsubo2016}, whose full profinite statement was
opened and checked. It does not use a power-series-domain argument.
The diagonal \((m,n)=(r,r)\) is cofinal in the divisibility product,
which proves the coefficient formulation in the statement.
\end{proof}

\begin{statement}{B.9}{Proved}{The exact permutation comparison}
On \(\Lambda c\) the selected \(S_3\)-action is semilinear:
\[
 r(f(X,Y)c)=f(Y,X^{-1}Y^{-1})c,\qquad
 s(f(X,Y)c)=-f(Y,X)c,                                       \tag{B.6}
\]
where \(r^3=s^2=(rs)^2=1\). It agrees with the coordinate-permutation
action on the Fermat tower after a compatible deck translation.
Changing that translation conjugates the action by the corresponding
deck operator. The same comparison is represented integrally at every
finite level by the matrices of B.7.
\end{statement}
\begin{proof}
Conjugation by the modular generators acts on \(F\) by
\(\rho:x\mapsto y,\ y\mapsto x^{-1}y^{-1}\) and
\(\sigma:x\mapsto y,\ y\mapsto x\).
Replace \(\rho\) by \(x\rho\); its based substitution is
\[
 r:x\mapsto xyx^{-1},\quad y\mapsto y^{-1}x^{-1}.
\]
Direct free reduction gives \(r(c)=c\) and \(s(c)=c^{-1}\).
Their actions on \(F^{\rm ab}\) are
\[
 R=\begin{pmatrix}0&-1\\1&-1\end{pmatrix},\qquad
 S=\begin{pmatrix}0&1\\1&0\end{pmatrix}.
\]
The defects in \(r^3\) and \((rs)^2\) are inner automorphisms by
elements of \(F'\). Such automorphisms act trivially on \(F'/F''\)
and on \(K_n^{\rm ab}\), since \(F'\subset K_n\). Formula (B.6)
therefore gives a genuine \(S_3\)-action.
At the level of deck extensions this is the complement constructed
in B.11. The coordinate permutations provide another complement
of the same \(A_n\)-extension. B.10 proves that there is a unique
deck translation conjugating the two. Uniqueness makes these
translations compatible as \(n\) varies. Passing to homology and
then to (B.5) gives the claimed comparison, with its ambiguity stated.
\end{proof}

\begin{statement}{B.10}{Proved}{The two cohomology groups governing uniqueness}
For every abelian group \(I\), with trivial \(S_3\)-action on \(I\),
put \(M_I=I^3/I(1,1,1)\). Then
\[
 H^0(S_3,M_I)=H^1(S_3,M_I)=0.                               \tag{B.7}
\]
The same assertions hold for continuous cohomology with a profinite
\(I\). For any profinite group \(B\) acting on \(I\) and commuting
with the permutation action,
\[
 H^0(B\times S_3,M_I)=H^1(B\times S_3,M_I)=0.
\]
\emph{Manuscript: \S51.12--13, folios 702, 715--716, scans 109,
136--138, (295)--(297); \S51.14, folio \(721'\), scan 148, (309).}
\end{statement}
\begin{proof}
Use coordinates in which \(r,s\) act by \(R,S\) above.
An invariant \((u,v)\) satisfies \(u=v\), \(u=-v\), and
\(u=2v\). Thus \(2v=3v=0\), so \(u=v=0\), over every \(I\).
A crossed homomorphism is specified by \(f(r)=a=(a_0,a_1)\),
\(f(s)=b\). The relation \(s^2=1\) gives \(b=(t,-t)\).
The relation \(r^3=1\) imposes no condition, because \(1+R+R^2=0\).
The relation \((rs)^2=1\) gives
\[
 a_0=2a_1+3t.
\]
The element \(m=(-a_1-2t,-a_1-t)\) satisfies
\((R-1)m=a\) and \((S-1)m=b\).
Hence the cocycle is a coboundary, and the conjugating element is
unique by the invariant calculation. Every formula uses only integral
addition, so it also proves the continuous profinite assertion.
For \(B\times S_3\), first subtract the unique coboundary that kills
the restriction to \(S_3\). The cocycle identity for commuting
\(b\in B,s\in S_3\) then says that \(f(b)\) is \(S_3\)-invariant.
It is zero. This proves the last assertion without a spectral-sequence
convergence assumption.
\end{proof}

\begin{statement}{B.11}{Proved}{An explicit integral splitting}
The extension \(\mathfrak S^+/\mu_{\mathfrak S}\) of \(S_3\) by \(M\)
splits. In the affine realization
\[
 \rho=(R,0),\qquad \sigma=(S,t),\qquad
 t=(1/3,-1/3),
\]
a complement is generated by
\[
 \rho'=\lambda_0\rho=(R,(1,0)),\qquad \sigma'=\sigma.
\]
It fixes \(p=(2/3,1/3)\). In particular its integral extension
class, and the profinite class after completion, are zero.
\emph{Manuscript: \S51.12, folios 700--702, scans 105--109;
\S51.13, folio \(715'\), scan 136.}
\end{statement}
\begin{proof}
For affine multiplication \((A,u)(B,v)=(AB,u+Av)\),
\(1+R+R^2=0\) gives \((\rho')^3=1\),
\((1+S)t=0\) gives \((\sigma')^2=1\), and
\[
 RS=\begin{pmatrix}-1&0\\-1&1\end{pmatrix},\qquad
 (1,0)+Rt=(4/3,2/3),\qquad
 (1+RS)(4/3,2/3)=0
\]
gives \((\rho'\sigma')^2=1\). Their linear parts generate all of
\(S_3\), so this is a section. Its translations lie in the original
affine cosets, since the only changed generator was multiplied by
\(\lambda_0\in M\). Solving \((1-R)p=(1,0)\) and
\((1-S)p=t\) gives the stated fixed point.

For complete cocycle data choose \(h=r^i s^j\) and put
\(t_h=R^i(jt)\). The original integral cocycle is
\[
 c(h,k)=t_h+A_h t_k-t_{hk}.
\]
The section just found has translation \(p-A_hp\).
Thus the integral cochain \(b_h=p-A_hp-t_h\) satisfies
\(c+\delta b=0\). This proves vanishing as a cohomology class,
not merely satisfaction of a few relations. All formulas commute
with reduction modulo \(N\) and with profinite completion.
\end{proof}

\begin{statement}{B.12}{Proved}{Splitting and uniqueness for the full quotient}
The continuous extension
\[
 1\longrightarrow\widehat M\longrightarrow
 G_{\rho,\sigma}/\mu_G\longrightarrow
 \widehat{\mathbb Z}^{\,*}\times S_3\longrightarrow1
\]
has zero class in continuous \(H^2\). Its set of continuous splittings
is a torsor under \(\widehat M\), by conjugation.
For the canonical scalar lift \(u_\mu\) of the manuscript, put
\(\nu=(\mu-1)/2\in\widehat{\mathbb Z}\). A section commuting with the
complement of B.11 is
\[
 s(\mu)=(-\nu,-\nu)\,u_\mu.
\]
\emph{Manuscript: \S51.12, folios 704--\(704'\), scans 113--114,
(264)--(266); \S51.13, folios 715--\(716'\), scans 135--138.}
\end{statement}
\begin{proof}
The manuscript's scalar automorphism sends
\(\rho\) to \(\lambda_1^\nu\rho\) and fixes \(\sigma\) modulo the
central involution, and it acts by multiplication by \(\mu\) on \(M\).
In the affine realization it is conjugation by
\[
 (\mu I,w_\mu),\qquad w_\mu=(-\nu/3,\nu/3).
\]
Indeed \((1-R)w_\mu=\nu(0,1)\) and
\(\mu t+(1-S)w_\mu=t\). The product law follows from
\(\nu(\mu\mu')=\nu(\mu)+\mu\nu(\mu')\).
After multiplying by the integral translation \((-\nu,-\nu)\),
its affine translation becomes \((1-\mu)p\).
These are precisely the scalar homotheties about \(p\); they commute
with the finite complement fixing \(p\), and multiply according to
\(\mu\mu'\). This gives a continuous section of the full product.
The thirds in this verification are computed in
\(\widehat{\mathbb Z}\otimes_{\mathbb Z}\mathbb Q\); the correcting
cochain is integral at every prime, including \(3\).
The half is integral because every profinite unit is odd at \(2\).

Two sections differ by a continuous \(1\)-cocycle with coefficients
\(\widehat M\), acted on by \((\mu,h)\) as \(\mu A_h\).
B.10 makes that cocycle a coboundary and makes the conjugating
element unique. Thus the asserted torsor and the vanishing of the
actual \(H^2\)-class both hold.
\end{proof}

\begin{statement}{B.13}{Proved}{The finite \(2\)- and \(3\)-primary checks}
For \(N\mid24\), the actual extension class of B.12 in coefficients
\(M/NM\) is zero, with trivial \(2\)- and \(3\)-primary parts.
The independent normalized-cochain computation gives
\[
 \left|H^0(S_3,M/NM)\right|=\left|H^1(S_3,M/NM)\right|=1,
 \qquad \left|H^2(S_3,M/NM)\right|=\gcd(N,3).
\]
In particular the zero answer for the actual class at \(N=3\)
is a genuine cocycle computation in a nonzero obstruction group.
\end{statement}
\begin{proof}
The integral cochain of B.11 is included in
\texttt{code/b\_splitting.json}. For the full product we use the
reference affine lift
\[
 t_{\mu,h}=w_\mu+\mu t_h,\qquad
 b_{\mu,h}=(-\nu,-\nu)+\mu b_h.
\]
The program verifies \(c+\delta b=0\) modulo \(N\) for every pair of
elements of \((\mathbb Z/6N)^*\times S_3\). The modulus \(6N\)
retains the halves and thirds in these reference formulas.
For \(N=1,2,3,4,6,8,12,24\) the numbers of checked pairs are,
respectively,
\[
 144,\ 576,\ 1296,\ 2304,\ 5184,\ 9216,\ 20736,\ 82944.
\]
The normalized differentials \(C^0\to C^1\to C^2\to C^3\)
are integer matrices. Their nonzero Smith diagonals are
\(1^2,\ 1^8,\ 1^{41},3\).
For an integer matrix with nonzero Smith entries \(d_i\), its kernel
modulo \(N\) has cardinality \(N^{b-r}\prod_i\gcd(N,d_i)\),
where \(b\) is its column count and \(r\) its rank.
Applying this formula to those three matrices gives the displayed
cohomology orders. The code checks consecutive differentials compose
to zero and verifies the full-product cocycle by exact rational
arithmetic before reduction.
\end{proof}

\begin{statement}{B.14}{Proved}{Identification and rigidity of the arithmetic tower}
Let
\[
 E=\pi_1^{\mathrm{et}}([U_1^+/S_3],\bar b),\qquad
 E_{\rm ab}=E/\overline{[\pi_1(U_{1,\overline{\mathbb Q}}),
                              \pi_1(U_{1,\overline{\mathbb Q}})]}.
\]
Its exact extension
\[
 1\to\widehat M(1)\to E_{\rm ab}\to\Gamma_{\mathbb Q}\times S_3\to1
\]
splits, with one conjugacy class of splittings and a unique
\(\widehat M(1)\)-conjugator between any two. The affine representation
on a Fermat fibre has target
\(\widehat M\rtimes(\widehat{\mathbb Z}^{\,*}\times S_3)\).
The representation supplied by the quotient
\(G_{\rho,\sigma}/\mu_G\), with the same geometric translation map
and cyclotomic character, agrees with it after a unique translation
of the fibre. This proves the identification posed in B.1.
\emph{Manuscript: \S51.13, folios 715--\(718'\), scans 135--142,
(294)--(298).}
\end{statement}
\begin{proof}
The coordinate-permutation action is defined over \(\mathbb Q\).
Consequently the arithmetic quotient of the orbifold group by the
geometric fundamental group is
\(\Gamma_{\mathbb Q}\times S_3\).
The fibres of \(U_n^+\to U_1^+\) form compatible torsors under
\(M/nM(1)\). Choose one point in their inverse-limit fibre.
Monodromy then acts by affine transformations with linear action
\(\chi(\gamma)A_h\), and its restriction to the geometric group is
the abelianization map of B.2. Thus it factors through \(E_{\rm ab}\).
The map from \(E_{\rm ab}\) to the pullback of this affine group over
\(\Gamma_{\mathbb Q}\times S_3\) is an isomorphism: it is the identity
on its kernel \(\widehat M(1)\) and on its quotient, so injectivity
and surjectivity follow by lifting a quotient element and correcting
in the kernel. The pullback of a semidirect product is split.

B.10 with \(B=\Gamma_{\mathbb Q}\), acting through \(\chi\), proves
uniqueness of the conjugacy class and of its conjugator.
The algebraically constructed representation and the Fermat
representation have the same restriction to \(\widehat M(1)\)
and the same quotient action. Their translation difference is
therefore a cocycle on \(\Gamma_{\mathbb Q}\times S_3\).
It is the unique coboundary given by B.10, which proves the
claimed comparison. This identifies the resulting unbased torsor
canonically; a based homomorphism still requires the fibre point
chosen above. The group \(E\) in this argument is the arithmetic
fundamental group, with quotient the actual Galois group.
No equality of that group with a group of all loop-preserving
outer automorphisms is used.

The source changes its tilde notation here. At scan 136 it describes
the geometric-abelian quotient, while at scans 139--140, in (298),
it uses \(\widetilde{\mathcal M}=\mathcal M/\mu_{\mathcal M}\)
for the orbifold group. Our notation retains the successive maps
\(\mathcal M\to E\to E_{\rm ab}\). Thus the argument uses the
displayed kernels, rather than identifying the two quotients
because they share a symbol in the manuscript.
\end{proof}

\paragraph{Scope of the two ranks.}
The map in B.14 first uses the quotient
\(\widehat F\to\widehat F^{\rm ab}=\widehat M(1)\).
The \(\Lambda\)-module in B.8 is obtained by taking the homology
\emph{of the kernels} of its finite quotients. This is the exact
passage from the universal abelian covering to the commutator
module seen by the metabelian fundamental group. It is this map
that connects the rank-two construction in \S51 to the universal
homology module.
